\documentclass{article}
\usepackage{amsmath}
\usepackage{amssymb}



\author{darthpedroo}
\title{Introducción a la programación resueltos}
\date{Segundo Cuatrimestre 2026}

\begin{document}
\maketitle

\section{Práctica 1}

\section{Práctica 2}

\textbf{Ejercicio 1.} Escribir semiformalmente los siguientes predicados

\begin{enumerate}
\item{\textit{esPrimo}: que dado un número verifica si cumple las propiedades de ser un número primo.}
\item{\textit{esPosicionValida}: que dado un entero i y una secuencia l, verifica si i es un índice válido para l.}
\item{\textit{esMinimo}: que dado una secuencia de enteros l y un entero \textit{elem},verifica que \textit{elem} sea el mínimo.}
\item{\textit{esMaximo}: que dada una secuencia de enteros \textit{l} y un entero \textit{elem}, verifica que \textit{elem} sea el máximo.}
\end{enumerate}

\textbf{Respuestas}

\[
\begin{array}{l}
\mathrm{proc}\ \mathit{esMinimo}\,(\mathrm{in}\ l : \mathit{seq}\langle\mathbb{Z}\rangle,\ \mathrm{in}\ \mathit{elem} : \mathbb{Z}) : \mathrm{Bool}\ \{ \\
\qquad \mathit{requiere}\ \{\mathrm{True}\} \\
\qquad \mathit{asegura}\ \{\mathit{res} = \mathrm{true} \leftrightarrow (\mathit{elem} \in l \wedge (\forall x : \mathbb{Z})(x \in l \rightarrow \mathit{elem} \leq x))\} \\
\}
\end{array}
\]

\end{document}
